\chapter{Evaluación experimental del sistema predictivo}\label{cap:resultados}

Este capítulo es el núcleo empírico del TFM: presenta y discute los resultados de los
once modelos descritos o derivados en el capítulo \ref{cap:metodologia}, evaluados sobre
las particiones de validación y test definidas en la sección \ref{sec:particiones_fuga}.
Ninguna cifra se introduce manualmente: todas llegan al documento desde los artefactos
persistidos por las fases 3 a 5b, por la cadena que detalla el Anexo \ref{cap:anexo3}.
La respuesta al contraste central del TFM (\ref{sec:contraste_central}) se
somete a tres controles diagnósticos independientes: la anatomía del residuo de la etapa 1
(\ref{sec:anatomia_residuo_etapa1}), la paridad informativa entre etapas
(\ref{sec:paridad_informativa}) y el desglose por subgrupos y periodos con corrección por
comparaciones múltiples (\ref{sec:desglose_subgrupos}).

\section{Tabla maestra de comparación}\label{sec:tabla_maestra}

Las tablas \ref{tab:comparacion_test} y \ref{tab:comparacion_val} reúnen los once modelos
de este TFM (cuatro baselines ingenuos, dos modelos de una sola etapa, tres variantes de
la familia LSTM y dos variantes de ensamble) ordenados de mejor a peor por MAE sobre las
particiones de test y de validación respectivamente. El coeficiente de determinación se
reporta a cuatro decimales y no a dos, una asimetría deliberada frente al resto de
métricas: a dos decimales, \texttt{ensemble\_equal} (0,9755), \texttt{xgboost\_alone}
(0,9690) y \texttt{sarimax} (0,9672) colapsarían en 0,98/0,97/0,97, borrando exactamente
la distinción que esta comparación existe para establecer. El MAPE, en cambio, no sufre
ese problema a esta escala y se mantiene a dos decimales.

\begin{table}[htbp]
  \centering
  \small
  \setlength{\tabcolsep}{4pt}
  \resizebox{\textwidth}{!}{\begin{tabular}{>{\raggedright\arraybackslash}p{3.2cm}>{\raggedright\arraybackslash}p{2.4cm}>{\centering\arraybackslash}p{0.8cm}>{\centering\arraybackslash}p{1.7cm}>{\centering\arraybackslash}p{1.7cm}>{\centering\arraybackslash}p{1.7cm}>{\centering\arraybackslash}p{1.4cm}}
\toprule
\rowcolor{naranja}
\textbf{Modelo} & \textbf{Familia} & \textbf{n} & \textbf{MAE} & \textbf{RMSE} & \textbf{MAPE (\%)} & \textbf{R²} \\
\midrule
\rowcolor{filaclara} ensemble\_\allowbreak{}inverse\_\allowbreak{}mae & Ensemble & 197 & 139.681 & 208.339 & 3,11 & 0,9755 \\
ensemble\_\allowbreak{}equal & Ensemble & 197 & 139.725 & 208.464 & 3,13 & 0,9755 \\
\rowcolor{filaclara} xgboost\_\allowbreak{}alone & Una etapa & 197 & 156.121 & 234.154 & 3,30 & 0,9690 \\
sarimax & Una etapa & 197 & 163.627 & 241.147 & 3,85 & 0,9672 \\
\rowcolor{filaclara} hybrid\_\allowbreak{}weighted & Familia LSTM & 197 & 220.286 & 319.447 & 5,17 & 0,9424 \\
hybrid & Familia LSTM & 197 & 234.223 & 328.131 & 5,49 & 0,9392 \\
\rowcolor{filaclara} lstm\_\allowbreak{}alone & Familia LSTM & 197 & 280.952 & 431.714 & 6,62 & 0,8948 \\
seasonal\_\allowbreak{}naive & Baseline ingenuo & 197 & 309.817 & 690.480 & 7,16 & 0,7308 \\
\rowcolor{filaclara} persistence & Baseline ingenuo & 197 & 900.570 & 1.390.444 & 21,79 & -0,0916 \\
moving\_\allowbreak{}average\_\allowbreak{}7 & Baseline ingenuo & 197 & 1.093.679 & 1.260.492 & 28,13 & 0,1029 \\
\rowcolor{filaclara} moving\_\allowbreak{}average\_\allowbreak{}28 & Baseline ingenuo & 197 & 1.119.987 & 1.301.939 & 28,74 & 0,0430 \\
\bottomrule
\end{tabular}
}
  \caption{\headlinecolor{Comparación de modelos, partición TEST (n=197, 2026-01-18 a 2026-08-02).}}
  \label{tab:comparacion_test}
\end{table}

\begin{table}[htbp]
  \centering
  \small
  \setlength{\tabcolsep}{4pt}
  \resizebox{\textwidth}{!}{\begin{tabular}{>{\raggedright\arraybackslash}p{3.2cm}>{\raggedright\arraybackslash}p{2.4cm}>{\centering\arraybackslash}p{0.8cm}>{\centering\arraybackslash}p{1.7cm}>{\centering\arraybackslash}p{1.7cm}>{\centering\arraybackslash}p{1.7cm}>{\centering\arraybackslash}p{1.4cm}}
\toprule
\rowcolor{naranja}
\textbf{Modelo} & \textbf{Familia} & \textbf{n} & \textbf{MAE} & \textbf{RMSE} & \textbf{MAPE (\%)} & \textbf{R²} \\
\midrule
\rowcolor{filaclara} ensemble\_\allowbreak{}inverse\_\allowbreak{}mae & Ensemble & 196 & 226.418 & 323.348 & 5,62 & 0,9519 \\
ensemble\_\allowbreak{}equal & Ensemble & 196 & 227.588 & 324.752 & 5,69 & 0,9515 \\
\rowcolor{filaclara} xgboost\_\allowbreak{}alone & Una etapa & 196 & 244.843 & 355.027 & 5,68 & 0,9420 \\
sarimax & Una etapa & 196 & 273.467 & 404.153 & 7,31 & 0,9248 \\
\rowcolor{filaclara} hybrid\_\allowbreak{}weighted & Familia LSTM & 196 & 276.347 & 401.456 & 6,94 & 0,9258 \\
hybrid & Familia LSTM & 196 & 280.738 & 415.562 & 7,09 & 0,9205 \\
\rowcolor{filaclara} lstm\_\allowbreak{}alone & Familia LSTM & 196 & 411.142 & 655.479 & 11,81 & 0,8023 \\
seasonal\_\allowbreak{}naive & Baseline ingenuo & 196 & 460.615 & 883.627 & 12,63 & 0,6408 \\
\rowcolor{filaclara} persistence & Baseline ingenuo & 196 & 914.119 & 1.387.393 & 23,72 & 0,1144 \\
moving\_\allowbreak{}average\_\allowbreak{}7 & Baseline ingenuo & 196 & 1.068.961 & 1.237.156 & 28,86 & 0,2958 \\
\rowcolor{filaclara} moving\_\allowbreak{}average\_\allowbreak{}28 & Baseline ingenuo & 196 & 1.137.759 & 1.357.281 & 31,12 & 0,1524 \\
\bottomrule
\end{tabular}
}
  \caption{\headlinecolor{Comparación de modelos, partición VAL (n=196, 2025-07-06 a 2026-01-17).}}
  \label{tab:comparacion_val}
\end{table}

La jerarquía resultante es nítida y consistente entre ambas particiones, como confirma
la figura \ref{fig:model_comparison_bar_test}.
En primer lugar, las dos variantes de ensamble (combinaciones puramente aritméticas de
SARIMAX y \texttt{xgboost\_alone}, sin entrenamiento adicional) lideran ambas
particiones: \texttt{ensemble\_inverse\_mae} obtiene el mejor MAE de test, \textbf{139.681}, con un
$R^2$ de 0,9755, seguido a menos de un 0,05\,\% de distancia por \texttt{ensemble\_equal}
(139.725). En segundo lugar, los dos modelos de una sola etapa ocupan el siguiente escalón
sin solaparse con el anterior: \texttt{xgboost\_alone} (MAE test 156.121, $R^2=0{,}9690$) y
SARIMAX (MAE test 163.627, $R^2=0{,}9672$) mejoran de forma clara a toda la familia LSTM
que seguirá describiéndose en las secciones siguientes, pero quedan por detrás del
ensamble en aproximadamente un 11\,\% y un 17\,\% de MAE respectivamente. En tercer
lugar, la familia LSTM (\texttt{hybrid\_weighted}, \texttt{hybrid} y \texttt{lstm\_alone},
en ese orden) se sitúa sistemáticamente peor que ambos modelos de una sola etapa, un
resultado que se analiza en detalle en la sección \ref{sec:contraste_central}.

Por último, los cuatro baselines ingenuos colapsan frente a cualquiera de los modelos
anteriores: la persistencia estacional, el mejor de ellos, supera en torno a un 10\,\% en
MAE al peor miembro de la familia LSTM, y la persistencia simple y las dos medias móviles presentan
$R^2$ próximos a cero o directamente negativos, confirmando sobre el conjunto completo de
modelos la advertencia ya adelantada en el capítulo \ref{cap:metodologia}: superar a la
persistencia simple es un umbral trivial en esta serie, y el resultado interesante de este
capítulo se juega enteramente entre los siete modelos restantes.

\begin{figure}[htbp]
  \centering
  \includegraphics[width=0.78\textwidth]{model_comparison_bar_test.png}
  \caption{\headlinecolor{Comparación de modelos en la partición test. Barras agrupadas de MAE, RMSE, MAPE y $R^2$, ordenadas de mejor a peor por MAE. El color codifica la familia del modelo.}}
  \label{fig:model_comparison_bar_test}
\end{figure}

Ninguna cifra de la partición de entrenamiento se presenta junto a estas dos tablas como
si fuera una estimación válida de desempeño: train es, por construcción, una muestra
\emph{in-sample} para todo modelo que se ajusta sobre ella, y su inclusión aquí induciría
a compararla incorrectamente con val o test, un error que la sección
\ref{sec:analisis_residuos} evita explícitamente al tratar el desempeño en train como
diagnóstico de sobreajuste y no como un resultado más de la tabla maestra.

\section{Ajuste temporal}\label{sec:ajuste_temporal}

La figura \ref{fig:demand_actual_vs_predicted_test} muestra la serie real de demanda
total superpuesta a la predicción del modelo ganador, \texttt{ensemble\_equal}, sobre la
partición de test; la de validación reproduce el mismo patrón. La predicción sigue de
cerca tanto la oscilación semanal (el diente de sierra laborable/fin de semana de la
sección \ref{sec:eda}) como las transiciones asociadas a festivos aislados, sin el retraso
de fase que delataría un modelo que en realidad predijera el valor del día anterior. Las
mayores desviaciones puntuales se concentran en los días de calendario irregular, festivos
y puentes, desglosados en la sección \ref{sec:error_dia}; el resto de la serie se sigue
con una fidelidad coherente con el MAPE de test del 3,13\,\% de la tabla
\ref{tab:comparacion_test}.

\begin{figure}[htbp]
  \centering
  \includegraphics[width=0.78\textwidth]{demand_actual_vs_predicted_test.png}
  \caption{\headlinecolor{Demanda diaria real y predicha por el modelo \texttt{ensemble\_equal} en la partición test. Serie observada y predicha superpuestas.}}
  \label{fig:demand_actual_vs_predicted_test}
\end{figure}

\section{Análisis de residuos}\label{sec:analisis_residuos}

El análisis de residuos permite distinguir un modelo que generaliza de uno que memoriza.
El caso más claro de esta distinción en todo el proyecto es \texttt{xgboost\_alone}: su
coeficiente de determinación en entrenamiento alcanza 0,9929, un ajuste casi perfecto,
frente a un MAE de test de 156.121; esa brecha entre un $R^2$ de entrenamiento
prácticamente unitario y un error de test sustancial es indicativa de que el modelo
memoriza una parte relevante de su conjunto de entrenamiento en lugar de limitarse a
capturar la estructura genuinamente generalizable de la demanda. La figura
\ref{fig:residual_distribution_xgboost_alone} hace visible esta brecha comparando las
distribuciones de error en train y en test: la primera es marcadamente más estrecha y
centrada que la segunda, un patrón de sobreajuste clásico en un modelo de árboles con
capacidad suficiente para ajustar el ruido de una muestra de apenas 889 observaciones
diarias. Es importante subrayar que la cifra que debe citarse en cualquier afirmación
sobre el desempeño de este modelo es siempre la de test, nunca la de entrenamiento como si
fuera indicativa de generalización.

\begin{figure}[htbp]
  \centering
  \includegraphics[width=0.78\textwidth]{residual_distribution_xgboost_alone.png}
  \caption{\headlinecolor{Distribución del error de predicción del modelo \texttt{xgboost\_alone} en la partición test. La brecha train-test es la evidencia visual del sobreajuste ($R^2$ train = 0,9929 frente a MAE test = 156.121).}}
  \label{fig:residual_distribution_xgboost_alone}
\end{figure}

Los modelos derivados de la etapa 1 LSTM, \texttt{hybrid} entre ellos, no disponen de esta
misma comparación train-test, puesto que la fase 4 únicamente persistió predicciones del
modelo final sobre las particiones de validación y de test, no sobre entrenamiento; no se
fuerza aquí una comparación que los datos disponibles no permiten. La
figura \ref{fig:residual_distribution_hybrid} muestra, en su lugar, la distribución del
error de \texttt{hybrid} en test sin más: una distribución visiblemente más dispersa que
la de \texttt{xgboost\_alone}, coherente con el MAE de test casi 50\,\% mayor que se
cuantificó en la sección \ref{sec:tabla_maestra}.

\begin{figure}[htbp]
  \centering
  \includegraphics[width=0.78\textwidth]{residual_distribution_hybrid.png}
  \caption{\headlinecolor{Distribución del error de predicción del modelo \texttt{hybrid} en la partición test.}}
  \label{fig:residual_distribution_hybrid}
\end{figure}

Por contraste, la figura \ref{fig:residual_distribution_ensemble_equal} muestra la
distribución de error del modelo ganador, \texttt{ensemble\_equal}: una distribución
sensiblemente más estrecha y mejor centrada en torno a cero que la de cualquiera de sus
dos componentes individuales por separado. Este estrechamiento es coherente con el
mecanismo de reducción de varianza por descorrelación de errores que se cuantifica en la
sección \ref{sec:mecanismo_ensemble}: promediar dos modelos cuyos errores no están
perfectamente correlacionados produce un error agregado con menor dispersión que
cualquiera de los dos errores originales, incluso sin que ninguno de los dos componentes
cambie en absoluto.

\begin{figure}[htbp]
  \centering
  \includegraphics[width=0.78\textwidth]{residual_distribution_ensemble_equal.png}
  \caption{\headlinecolor{Distribución del error de predicción del modelo \texttt{ensemble\_equal} (ganador) en la partición test.}}
  \label{fig:residual_distribution_ensemble_equal}
\end{figure}

\section{Error por tipo de día}\label{sec:error_dia}

Las tablas \ref{tab:desglose_dia_test} y \ref{tab:desglose_dia_val} descomponen el MAE de
cada modelo por régimen de calendario (laborable ordinario, laborable con puente, sábado,
domingo y festivo), y la figura \ref{fig:error_by_day_type_test} lo representa
gráficamente, atenuando visualmente los
grupos con menos de diez observaciones y anotando su tamaño muestral bajo cada barra. La
sección \ref{sec:desglose_subgrupos} retoma esta descomposición de forma sistemática
(quince subgrupos de calendario y periodo fijados por adelantado, con réplica en validación
y corrección por comparaciones múltiples) para responder a si el híbrido mejora en algún
régimen concreto.

\begin{table}[htbp]
  \centering
  \small
  \begin{tabular}{>{\raggedright\arraybackslash}p{2.6cm}>{\centering\arraybackslash}p{0.7cm}>{\centering\arraybackslash}p{1.7cm}>{\centering\arraybackslash}p{1.9cm}>{\centering\arraybackslash}p{1.6cm}>{\centering\arraybackslash}p{1.9cm}>{\centering\arraybackslash}p{1.9cm}}
\toprule
\rowcolor{naranja}
\textbf{Tipo de día} & \textbf{n} & \textbf{sarimax} & \textbf{xgboost\_\allowbreak{}alone} & \textbf{hybrid} & \textbf{hybrid\_\allowbreak{}weighted} & \textbf{ensemble\_\allowbreak{}equal} \\
\midrule
\rowcolor{filaclara} Laborable ordinario & 133 & 148.297 & 169.488 & 211.581 & 200.326 & 138.628 \\
Laborable & 136 & 150.493 & 169.845 & 215.729 & 205.438 & 140.352 \\
\rowcolor{filaclara} Sábado & 27 & 203.910 & 129.381 & 247.793 & 185.918 & 138.805 \\
Domingo & 29 & 169.129 & 100.200 & 233.778 & 243.096 & 122.591 \\
\rowcolor{filaclara} Festivo & 5 & 271.410 & 251.553 & 666.572 & 677.429 & 227.037 \\
Laborable, puente & 3 & 247.865 & 185.671 & 399.637 & 432.066 & 216.768 \\
\bottomrule
\end{tabular}

  \caption{\headlinecolor{MAE por tipo de día, partición test. n=5 (festivo) y n=3 (laborable, puente): indicativos, no concluyentes. La fila ``Laborable'' ($n=136$) es la suma de ``Laborable ordinario'' (133) y ``Laborable, puente'' (3); no es un grupo adicional independiente.}}
  \label{tab:desglose_dia_test}
\end{table}

\begin{table}[htbp]
  \centering
  \small
  \begin{tabular}{>{\raggedright\arraybackslash}p{2.6cm}>{\centering\arraybackslash}p{0.7cm}>{\centering\arraybackslash}p{1.7cm}>{\centering\arraybackslash}p{1.9cm}>{\centering\arraybackslash}p{1.6cm}>{\centering\arraybackslash}p{1.9cm}>{\centering\arraybackslash}p{1.9cm}}
\toprule
\rowcolor{naranja}
\textbf{Tipo de día} & \textbf{n} & \textbf{sarimax} & \textbf{xgboost\_\allowbreak{}alone} & \textbf{hybrid} & \textbf{hybrid\_\allowbreak{}weighted} & \textbf{ensemble\_\allowbreak{}equal} \\
\midrule
\rowcolor{filaclara} Laborable ordinario & 121 & 211.578 & 263.817 & 242.993 & 240.264 & 205.597 \\
Laborable & 133 & 238.197 & 285.078 & 286.849 & 280.851 & 231.011 \\
\rowcolor{filaclara} Sábado & 26 & 435.186 & 132.334 & 278.593 & 245.120 & 246.171 \\
Domingo & 28 & 245.694 & 128.658 & 245.450 & 253.547 & 162.840 \\
\rowcolor{filaclara} Festivo & 9 & 413.907 & 336.750 & 306.409 & 370.935 & 324.760 \\
Laborable, puente & 12 & 506.603 & 499.452 & 729.068 & 690.105 & 487.269 \\
\bottomrule
\end{tabular}

  \caption{\headlinecolor{MAE por tipo de día, partición val. La fila ``Laborable'' ($n=133$) es la suma de ``Laborable ordinario'' (121) y ``Laborable, puente'' (12); no es un grupo adicional independiente.}}
  \label{tab:desglose_dia_val}
\end{table}

Esta descomposición revela la razón mecánica por la que el ensamble supera a ambos
modelos de una sola etapa que lo componen, y no solo el hecho de que lo hace. Sobre la
partición test, SARIMAX es claramente mejor que \texttt{xgboost\_alone} en los días
laborables ordinarios (MAE 148.297 frente a 169.488): su componente autorregresivo
estacional captura con precisión la dinámica de la inmensa mayoría de los días del año,
que son precisamente laborables ordinarios. \texttt{xgboost\_alone}, en cambio, es
netamente superior en fin de semana: en domingo su MAE es 100.200 frente a los 169.129 de
SARIMAX, y en sábado 129.381 frente a 203.910, casi 40\,\% y 37\,\% mejor,
respectivamente. \textbf{Ningún modelo domina al otro de forma global}; cada uno acierta mejor en
un subconjunto distinto y complementario de los días del año, y es exactamente esa
complementariedad la que hace mecánicamente ventajoso promediarlos, como se formaliza en
la sección \ref{sec:mecanismo_ensemble}.

El híbrido y su variante reponderada no participan de esta complementariedad: son, con la
única excepción de la comparación frente a \texttt{lstm\_alone}, peores que SARIMAX y que
\texttt{xgboost\_alone} en prácticamente todos los grupos de la tabla
\ref{tab:desglose_dia_test}. El hallazgo parcialmente positivo, y el único genuinamente
favorable a la arquitectura híbrida en todo ese desglose, aparece en festivos: la etapa 2
reduce el error de su propia etapa 1 de 1.585.388 (fila festivo, columna
\texttt{lstm\_alone} de la tabla \ref{tab:desglose_dia_test}) a 666.572, una mejora del
58\,\%, que muestra de nuevo que
el mecanismo de corrección funciona como está diseñado. Aun así, ese 666.572 sigue siendo
2,5 veces peor que el MAE de SARIMAX en el mismo grupo (271.410) y notablemente peor que
\texttt{xgboost\_alone} (251.553): la corrección funciona, pero parte de una etapa 1
demasiado débil como para que el resultado corregido alcance a los modelos de una sola
etapa.

\begin{figure}[htbp]
  \centering
  \includegraphics[width=0.78\textwidth]{error_by_day_type_test.png}
  \caption{\headlinecolor{MAE por tipo de día en la partición test. Los grupos con $n < 10$ se dibujan con opacidad reducida y su n anotado bajo el eje: son indicativos, no concluyentes.}}
  \label{fig:error_by_day_type_test}
\end{figure}

Es indispensable acompañar toda esta lectura de una nota de cautela obligatoria: en la
partición test, el grupo de festivos contiene únicamente cinco observaciones y el de
laborables con puente solo tres. Con muestras de ese tamaño, un único día atípico puede
desplazar el MAE del grupo entero en cientos de miles de viajes, de modo que estas dos
filas deben leerse como indicativas y en ningún caso como concluyentes; para afirmaciones
robustas sobre \emph{dónde} se concentra estructuralmente el error de un modelo univariante
frente al calendario, la base más sólida sigue siendo el diagnóstico exploratorio de la
fase 4 sobre el periodo completo (48 festivos y 54 puentes; la codificación de
calendario detecta 55 días puente sobre el total de 1.310 días, de los que 54 quedan
fuera de los 28 días de calentamiento excluidos del entrenamiento), no las particiones
reducidas de test y validación empleadas para esta comparación final.

\section{Importancia de variables}\label{sec:importancia_variables}

Las figuras \ref{fig:feature_importance_group_xgboost_alone} y
\ref{fig:feature_importance_group_xgboost_residual}, junto con las tablas
\ref{tab:importancia_grupo_alone} y \ref{tab:importancia_grupo_residual}, agregan la
ganancia de cada uno de los dos modelos XGBoost del proyecto por grupo temático de
variable, y las tablas \ref{tab:top15_alone} y \ref{tab:top15_residual} descienden al
detalle de las quince variables individuales más influyentes de cada modelo. El contraste entre ambos modelos es la evidencia más directa
de todo el capítulo sobre qué tipo de información resuelve cada problema de predicción.

\begin{figure}[htbp]
  \centering
  \includegraphics[width=0.78\textwidth]{feature_importance_group_xgboost_alone.png}
  \caption{\headlinecolor{Importancia agregada por grupo de variables del modelo \texttt{xgboost\_alone}, medida como ganancia.}}
  \label{fig:feature_importance_group_xgboost_alone}
\end{figure}

En \texttt{xgboost\_alone}, el grupo de calendario, solo 11 de las 161 variables
disponibles, concentra el 58,5\,\% de la ganancia total del modelo (tabla
\ref{tab:importancia_grupo_alone}), y su variable individual más influyente con
diferencia, \texttt{feat\_day\_type\_laborable}, aporta por sí sola una ganancia de
0,561 (tabla \ref{tab:top15_alone}): al predecir la demanda directamente, saber si el día
es laborable resuelve la práctica totalidad de la varianza semanal identificada en la
sección \ref{sec:eda}, y el resto de variables (retardos de operador, meteorología
retardada, términos de Fourier) se reparten un papel claramente secundario de matiz fino
sobre ese nivel base.

\begin{table}[htbp]
  \centering
  \small
  \begin{tabular}{lrrr}
\toprule
\rowcolor{naranja}
\textbf{Grupo} & \textbf{N.º variables} & \textbf{Ganancia} & \textbf{Peso (\%)} \\
\midrule
\rowcolor{filaclara} calendar & 11 & 0,5846 & 58,46 \\
rolling & 4 & 0,1409 & 14,09 \\
\rowcolor{filaclara} weather & 105 & 0,1081 & 10,81 \\
lag\_\allowbreak{}operator & 28 & 0,1007 & 10,07 \\
\rowcolor{filaclara} lag\_\allowbreak{}total & 7 & 0,0375 & 3,75 \\
fourier\_\allowbreak{}annual & 4 & 0,0268 & 2,68 \\
\rowcolor{filaclara} fourier\_\allowbreak{}weekly & 2 & 0,0014 & 0,14 \\
\bottomrule
\end{tabular}

  \caption{\headlinecolor{Importancia por grupo de variables, \texttt{xgboost\_alone}.}}
  \label{tab:importancia_grupo_alone}
\end{table}

\begin{table}[htbp]
  \centering
  \small
  \begin{tabular}{rllr}
\toprule
\rowcolor{naranja}
\textbf{\#} & \textbf{Variable} & \textbf{Grupo} & \textbf{Ganancia} \\
\midrule
\rowcolor{filaclara} 1 & feat\_\allowbreak{}day\_\allowbreak{}type\_\allowbreak{}laborable & calendar & 0,5609 \\
2 & feat\_\allowbreak{}total\_\allowbreak{}roll\_\allowbreak{}std\_\allowbreak{}7 & rolling & 0,1321 \\
\rowcolor{filaclara} 3 & feat\_\allowbreak{}carretera\_\allowbreak{}lag\_\allowbreak{}1 & lag\_\allowbreak{}operator & 0,0209 \\
4 & feat\_\allowbreak{}total\_\allowbreak{}lag\_\allowbreak{}1 & lag\_\allowbreak{}total & 0,0157 \\
\rowcolor{filaclara} 5 & feat\_\allowbreak{}day\_\allowbreak{}type\_\allowbreak{}sabado & calendar & 0,0156 \\
6 & feat\_\allowbreak{}metro\_\allowbreak{}lag\_\allowbreak{}1 & lag\_\allowbreak{}operator & 0,0152 \\
\rowcolor{filaclara} 7 & feat\_\allowbreak{}metro\_\allowbreak{}lag\_\allowbreak{}14 & lag\_\allowbreak{}operator & 0,0134 \\
8 & feat\_\allowbreak{}fourier\_\allowbreak{}annual\_\allowbreak{}k2\_\allowbreak{}sin & fourier\_\allowbreak{}annual & 0,0127 \\
\rowcolor{filaclara} 9 & feat\_\allowbreak{}relative\_\allowbreak{}humidity\_\allowbreak{}2m\_\allowbreak{}mean\_\allowbreak{}lag\_\allowbreak{}2 & weather & 0,0111 \\
10 & feat\_\allowbreak{}emt\_\allowbreak{}lag\_\allowbreak{}1 & lag\_\allowbreak{}operator & 0,0110 \\
\rowcolor{filaclara} 11 & feat\_\allowbreak{}apparent\_\allowbreak{}temperature\_\allowbreak{}min\_\allowbreak{}lag\_\allowbreak{}1 & weather & 0,0086 \\
12 & feat\_\allowbreak{}total\_\allowbreak{}lag\_\allowbreak{}3 & lag\_\allowbreak{}total & 0,0075 \\
\rowcolor{filaclara} 13 & feat\_\allowbreak{}is\_\allowbreak{}bridge\_\allowbreak{}day & calendar & 0,0070 \\
14 & feat\_\allowbreak{}relative\_\allowbreak{}humidity\_\allowbreak{}2m\_\allowbreak{}mean\_\allowbreak{}roll\_\allowbreak{}mean\_\allowbreak{}7 & weather & 0,0064 \\
\rowcolor{filaclara} 15 & feat\_\allowbreak{}fourier\_\allowbreak{}annual\_\allowbreak{}k1\_\allowbreak{}cos & fourier\_\allowbreak{}annual & 0,0060 \\
\bottomrule
\end{tabular}

  \caption{\headlinecolor{Top-15 variables individuales por ganancia, \texttt{xgboost\_alone}.}}
  \label{tab:top15_alone}
\end{table}

El modelo residual de la etapa 2, \texttt{xgboost\_residual}, invierte por completo esta
jerarquía de grupos (figura \ref{fig:feature_importance_group_xgboost_residual}, tabla
\ref{tab:importancia_grupo_residual}): la meteorología retardada, con sus 105 variables,
concentra ahora el 51,6\,\% de la ganancia, mientras que el calendario desciende al
17,3\,\%, prácticamente empatado con los retardos de operador (17,2\,\%). La lectura
inmediata de esta inversión, y la que se adoptó al observarla por primera vez, es
coherente con la propia definición de la variable objetivo de esta etapa: el residuo
$e_t = y_t - \hat y^{\text{LSTM}}_t$ ya no contendría la componente de calendario más
gruesa, la distinción laborable/fin de semana, porque una parte relevante de esa
componente ya habría sido absorbida, con mayor o menor acierto, por la propia dinámica
temporal de la etapa 1; lo que quedaría por explicar en el residuo sería, en consecuencia,
un matiz más fino, y ese matiz estaría más ligado a las condiciones meteorológicas
retardadas que al calendario grueso.

Conviene subrayar que esto es una \emph{hipótesis}
sugerida por una clasificación de importancia, no un resultado medido, y que las fases
diagnósticas la someten a prueba con dos análisis independientes y la \textbf{rechazan}:
la sección \ref{sec:residuo_meteo_anatomia} no encuentra ninguna correlación parcial
significativa entre el residuo y las 105 variables meteorológicas retardadas, y la
ablación de la sección \ref{sec:ablacion_bloques} muestra que retirar ese bloque
\emph{mejora} el modelo. La sección \ref{sec:lectura_conjunta} reconcilia esa aparente
contradicción y explica por qué una cuota de ganancia elevada es perfectamente compatible
con la ausencia de señal aprovechable.

\begin{figure}[htbp]
  \centering
  \includegraphics[width=0.78\textwidth]{feature_importance_group_xgboost_residual.png}
  \caption{\headlinecolor{Importancia agregada por grupo de variables del modelo \texttt{xgboost\_\allowbreak{}residual} (etapa 2 del híbrido).}}
  \label{fig:feature_importance_group_xgboost_residual}
\end{figure}

\begin{table}[htbp]
  \centering
  \small
  \begin{tabular}{lrrr}
\toprule
\rowcolor{naranja}
\textbf{Grupo} & \textbf{N.º variables} & \textbf{Ganancia} & \textbf{Peso (\%)} \\
\midrule
\rowcolor{filaclara} weather & 105 & 0,5157 & 51,57 \\
calendar & 11 & 0,1726 & 17,26 \\
\rowcolor{filaclara} lag\_\allowbreak{}operator & 28 & 0,1723 & 17,23 \\
lag\_\allowbreak{}total & 7 & 0,0592 & 5,92 \\
\rowcolor{filaclara} rolling & 4 & 0,0405 & 4,05 \\
fourier\_\allowbreak{}annual & 4 & 0,0320 & 3,20 \\
\rowcolor{filaclara} fourier\_\allowbreak{}weekly & 2 & 0,0077 & 0,77 \\
\bottomrule
\end{tabular}

  \caption{\headlinecolor{Importancia por grupo de variables, \texttt{xgboost\_residual}.}}
  \label{tab:importancia_grupo_residual}
\end{table}

Sin embargo, esta relectura por grupo no contradice el hallazgo mecanístico más
importante de la etapa 2, visible únicamente al bajar al nivel de variable individual
(tabla \ref{tab:top15_residual}): pese a que el calendario pesa menos como grupo agregado,
la variable individual con mayor ganancia en todo el modelo residual sigue siendo
\texttt{feat\_day\_type\_festivo}, con una ganancia de 0,125, muy por delante de la
segunda variable de la lista. Esto es coherente, al nivel de variable individual, con que la
etapa 2 aprenda la señal de calendario (en concreto, la irregularidad de los festivos) que
la etapa 1 univariante no puede ver por construcción: el mecanismo que justifica
teóricamente la arquitectura híbrida se manifiesta en la variable que cabría esperar, con
independencia de que el veredicto agregado de la
sección \ref{sec:contraste_central} termine siendo negativo para la arquitectura completa.

\begin{table}[htbp]
  \centering
  \small
  \begin{tabular}{rllr}
\toprule
\rowcolor{naranja}
\textbf{\#} & \textbf{Variable} & \textbf{Grupo} & \textbf{Ganancia} \\
\midrule
\rowcolor{filaclara} 1 & feat\_\allowbreak{}day\_\allowbreak{}type\_\allowbreak{}festivo & calendar & 0,1254 \\
2 & feat\_\allowbreak{}apparent\_\allowbreak{}temperature\_\allowbreak{}max\_\allowbreak{}lag\_\allowbreak{}7 & weather & 0,0759 \\
\rowcolor{filaclara} 3 & feat\_\allowbreak{}temperature\_\allowbreak{}2m\_\allowbreak{}mean\_\allowbreak{}lag\_\allowbreak{}3 & weather & 0,0337 \\
4 & feat\_\allowbreak{}carretera\_\allowbreak{}lag\_\allowbreak{}3 & lag\_\allowbreak{}operator & 0,0319 \\
\rowcolor{filaclara} 5 & feat\_\allowbreak{}day\_\allowbreak{}type\_\allowbreak{}laborable & calendar & 0,0305 \\
6 & feat\_\allowbreak{}total\_\allowbreak{}lag\_\allowbreak{}7 & lag\_\allowbreak{}total & 0,0249 \\
\rowcolor{filaclara} 7 & feat\_\allowbreak{}total\_\allowbreak{}roll\_\allowbreak{}std\_\allowbreak{}28 & rolling & 0,0217 \\
8 & feat\_\allowbreak{}apparent\_\allowbreak{}temperature\_\allowbreak{}min\_\allowbreak{}lag\_\allowbreak{}3 & weather & 0,0200 \\
\rowcolor{filaclara} 9 & feat\_\allowbreak{}apparent\_\allowbreak{}temperature\_\allowbreak{}max\_\allowbreak{}lag\_\allowbreak{}3 & weather & 0,0191 \\
10 & feat\_\allowbreak{}temperature\_\allowbreak{}2m\_\allowbreak{}min\_\allowbreak{}roll\_\allowbreak{}mean\_\allowbreak{}7 & weather & 0,0189 \\
\rowcolor{filaclara} 11 & feat\_\allowbreak{}weather\_\allowbreak{}code\_\allowbreak{}roll\_\allowbreak{}mean\_\allowbreak{}7 & weather & 0,0152 \\
12 & feat\_\allowbreak{}metro\_\allowbreak{}lag\_\allowbreak{}7 & lag\_\allowbreak{}operator & 0,0152 \\
\rowcolor{filaclara} 13 & feat\_\allowbreak{}shortwave\_\allowbreak{}radiation\_\allowbreak{}sum\_\allowbreak{}roll\_\allowbreak{}mean\_\allowbreak{}7 & weather & 0,0149 \\
14 & feat\_\allowbreak{}total\_\allowbreak{}lag\_\allowbreak{}28 & lag\_\allowbreak{}total & 0,0137 \\
\rowcolor{filaclara} 15 & feat\_\allowbreak{}temperature\_\allowbreak{}2m\_\allowbreak{}max\_\allowbreak{}roll\_\allowbreak{}mean\_\allowbreak{}7 & weather & 0,0128 \\
\bottomrule
\end{tabular}

  \caption{\headlinecolor{Top-15 variables individuales por ganancia, \texttt{xgboost\_residual}. \texttt{feat\_\allowbreak{}day\_\allowbreak{}type\_\allowbreak{}festivo} encabeza con ganancia 0,125 -- coherente con que la etapa 2 aprenda la señal de calendario que la etapa 1 no puede ver.}}
  \label{tab:top15_residual}
\end{table}

\section{Contraste central del TFM}\label{sec:contraste_central}

Con la evidencia de las secciones anteriores ya establecida, esta sección responde de
forma explícita a la pregunta que estructura todo el proyecto: ¿mejora la
arquitectura híbrida residual LSTM-XGBoost, tal como se ha diseñado e implementado en
este TFM, la predicción de demanda frente a SARIMAX y frente a un XGBoost entrenado
directamente sobre las mismas variables predictoras? La respuesta, sostenida por la tabla
\ref{tab:delta_hibrido}, es \textbf{no, en ninguna de las dos comparaciones y en ninguna
de las dos particiones de evaluación}.

\begin{table}[htbp]
  \centering
  \begin{tabular}{llrl}
\toprule
\rowcolor{naranja}
\textbf{Partición} & \textbf{Comparación} & \textbf{Delta MAE} & \textbf{Veredicto} \\
\midrule
\rowcolor{filaclara} val & hybrid - sarimax & 7.270 & hybrid peor \\
val & hybrid - xgboost\_\allowbreak{}alone & 35.895 & hybrid peor \\
\rowcolor{filaclara} test & hybrid - sarimax & 70.597 & hybrid peor \\
test & hybrid - xgboost\_\allowbreak{}alone & 78.103 & hybrid peor \\
\bottomrule
\end{tabular}

  \caption{\headlinecolor{Diferencia de MAE del híbrido frente a SARIMAX y XGBoost independiente (Delta MAE = \texttt{hybrid} $-$ competidor; positivo = híbrido peor). Generada, no introducida manualmente.}}
  \label{tab:delta_hibrido}
\end{table}

Sobre la partición test, el híbrido supera en MAE a SARIMAX en 70.597 viajes diarios y a
\texttt{xgboost\_alone} en 78.103; sobre validación las diferencias son menores en
magnitud absoluta, pero mantienen el mismo signo, 7.270 y 35.895 respectivamente. La
diferencia frente a \texttt{xgboost\_alone} es, en ambas particiones, mayor que la
diferencia frente a SARIMAX: el híbrido no solo pierde frente al baseline estadístico
clásico, sino que pierde por un margen todavía más amplio frente a un modelo de
aprendizaje automático que utiliza exactamente la misma matriz de variables predictoras
sin pasar por ninguna etapa LSTM. Este segundo resultado es el más informativo de los dos: si
el problema fuera únicamente que SARIMAX es un modelo particularmente bien ajustado a esta
serie, cabría esperar que el híbrido, con acceso a variables no lineales y a una capa de
aprendizaje profundo, cerrara al menos parte de esa distancia frente a otro modelo de
aprendizaje automático; en su lugar, la distancia se amplía.

\subsection{Por qué pierde el híbrido: lectura a la luz de Granger (1989)}\label{sec:resultado_negativo_teoria}

La teoría clásica de combinación de pronósticos ofrece un marco preciso para interpretar
este resultado, en lugar de limitarse a documentarlo. \textcite{granger1989combining},
revisando dos décadas de literatura sobre combinación de modelos, formula la condición que
determina cuándo un pronóstico híbrido o combinado puede superar a sus componentes: la
mejora solo es posible cuando los modelos que se combinan son individualmente subóptimos
\emph{y}, de manera crucial, se basan en conjuntos de información distintos entre sí. La
consecuencia de esa condición es que la ganancia de combinar procede de información que
cada componente tiene y el otro no; si el conjunto de información de uno está contenido en
el del otro, la combinación no puede aportar, en teoría, nada que un único modelo sobre el
conjunto mayor no pudiera aprender por sí mismo con una forma funcional suficientemente
flexible.

La arquitectura híbrida de este TFM \textbf{incumple precisamente esa segunda condición}, y
conviene precisar en qué sentido. Los conjuntos de información de las dos etapas no son
idénticos: la etapa 1 observa únicamente la ventana cruda de 28 valores de \texttt{total},
y la etapa 2 dispone además del calendario, la meteorología retardada, los armónicos de
Fourier y los retardos de operador —información que la etapa 1 no ve y que, como muestra
la sección \ref{sec:importancia_variables}, aporta señal genuina; de ahí que la etapa 2
mejore de forma sustancial a su propia etapa 1 (sección \ref{sec:error_dia})—. Pero
tampoco son \emph{distintos} en el sentido que la condición exige: el pasado reciente de la
serie que la etapa 1 observa está representado en la matriz de la etapa 2 por sus siete
retardos y cuatro estadísticos móviles sobre la misma ventana de 28 días, de modo que los
dos conjuntos están, en la práctica, anidados y no son complementarios. En esa situación la
combinación secuencial no puede superar, en teoría, al mejor modelo de una sola etapa sobre
el conjunto mayor, y ese modelo existe en la comparación: \texttt{xgboost\_alone}, que
accede a las mismas 161 variables predictoras sin arrastrar el ruido de una etapa 1
estructuralmente débil. Aunque la segunda etapa incorpora información adicional de
calendario, meteorología retardada y otras variables predictoras, los resultados indican
que esta separación secuencial no genera una complementariedad predictiva suficiente para
compensar el error introducido por la primera etapa LSTM. Los dos controles de la sección
\ref{sec:paridad_informativa} miden ambas mitades de esta lectura: la ablación muestra que
\texttt{xgboost\_alone} extrae por sí mismo la información temporal que la etapa 1 aporta,
y el control de paridad muestra que la vía recurrente no ofrece una representación de ese
pasado que el modelo de una sola etapa no tenga ya.

Esta lectura se refuerza, además, por un factor puramente cuantitativo: la etapa 1
univariante deja un residuo con una desviación típica de 726.669 viajes diarios
(\ref{sec:etapa2_xgboost}), del mismo orden de magnitud que la propia señal que se intenta
predecir. Cuanto mayor es la varianza residual que la segunda etapa debe corregir, más
difícil resulta que esa corrección alcance, y menos aún supere, a un modelo que ataca el
problema completo desde el principio con toda la matriz predictora disponible. Este patrón
—hibridar puede mejorar sustancialmente a un componente débil sin llegar a igualar a un
componente fuerte de una sola etapa cuando sus conjuntos de información están anidados— no
es exclusivo de
este proyecto: es consistente con lo observado de forma sistemática en las competiciones
de pronóstico a gran escala M4 y M5 \parencite{makridakis2018m4,makridakis2022m5}, donde los métodos híbridos y
de combinación que finalmente destacan lo hacen precisamente al combinar familias de
modelos con mecanismos e información de entrada genuinamente distintos, como la propia
sección \ref{sec:mecanismo_ensemble} muestra para el ensamble de este TFM, y no al
encadenar dos etapas sucesivas sobre conjuntos de información anidados. El resultado
negativo del híbrido residual de este TFM debe leerse, por tanto, como consistente con la
teoría publicada y con la evidencia empírica a gran escala, no como una anomalía
específica de esta implementación.

\section{Análisis de sensibilidad}\label{sec:analisis_sensibilidad}

Las dos secciones siguientes documentan dos experimentos acotados y registrados por
adelantado, con su criterio de éxito fijado antes de observar el resultado, diseñados
para explorar el margen de mejora del híbrido y para establecer un punto de comparación
adicional, respectivamente. Ninguno de los dos reabre la conclusión central de la sección
anterior ni constituye una búsqueda de una configuración que la revierta.

\subsection{Experimento A: repesado de días irregulares (3x)}\label{sec:experimento_a}

El primer experimento pregunta si concentrar la capacidad de ajuste de la etapa 2
precisamente en los días de calendario irregular —festivos y puentes, el subconjunto
donde la sección \ref{sec:error_dia} documentó el error relativo más alto de todo el
sistema— mejora el desempeño del híbrido en esos mismos días, a cambio de un coste
acotado en el desempeño global. La variante \texttt{hybrid\_weighted}, descrita en la
sección \ref{sec:modelos_control}, asigna un peso de muestra triple a esas filas (un
8,2\,\% del conjunto residual) sin modificar ningún otro elemento del modelo. El criterio
de éxito, fijado antes de observar el resultado y evaluado mecánicamente en código, exige
el cumplimiento simultáneo de dos condiciones: que el MAE en festivos y puentes mejore, y
que el MAE global de test no empeore más de un 5\,\% frente al híbrido original de la
sección \ref{sec:contraste_central}.

\begin{table}[htbp]
  \centering
  \setlength{\tabcolsep}{4pt}
  \resizebox{\textwidth}{!}{\begin{tabular}{>{\raggedright\arraybackslash}p{6.5cm}>{\centering\arraybackslash}p{1.9cm}>{\centering\arraybackslash}p{2.4cm}>{\centering\arraybackslash}p{3.0cm}}
\toprule
\rowcolor{naranja}
\textbf{Criterio} & \textbf{hybrid} & \textbf{hybrid\_\allowbreak{}weighted} & \textbf{Resultado} \\
\midrule
\rowcolor{filaclara} MAE en días irregulares (festivo + puente), n-ponderado & 566.472 & 585.418 & No mejora \\
MAE global test (techo +5\% sobre 234.223) & 234.223 & 220.286 & Dentro del margen \\
\rowcolor{filaclara} Veredicto del criterio prerregistrado (ambas condiciones) & -- & -- & NO CUMPLIDO \\
\bottomrule
\end{tabular}
}
  \caption{\headlinecolor{Veredicto del criterio prerregistrado del experimento A (repesado 3x de días irregulares). CRITERIO NO CUMPLIDO: el presupuesto global se cumplió pero el error en días irregulares NO mejoró.}}
  \label{tab:criterio_experimento_a}
\end{table}

El veredicto de la tabla \ref{tab:criterio_experimento_a} es que el criterio \textbf{no se
cumple}: la condición de presupuesto global sí se satisface, con un MAE de test que
incluso mejora un 5,95\,\% respecto al híbrido original (234.223 a 220.286), pero el MAE
ponderado por tamaño de muestra en festivos y puentes conjuntamente \emph{empeora}, de
566.472 a 585.418. El aspecto más interesante de este resultado nulo no es que el
repesado fallara, sino \emph{dónde} recayó la mejora que sí se produjo: no en los días que
se pretendía corregir, sino en sábados (mejora del 25\,\%) y en laborables ordinarios
(mejora del 5\,\%). Con solamente 45 filas repesadas frente a 161 variables de entrada, el
peso adicional no parece producir un mejor ajuste específicamente en esas filas; en su
lugar, perturba el ajuste global del modelo, y esa perturbación resulta beneficiar, por
las características propias del conjunto de datos, a los grupos más numerosos en lugar de
a los grupos objetivo. Se reporta como lo que es: un dato sobre la fragilidad del
repesado en muestras pequeñas, no como evidencia a favor ni en contra del esquema de
pesos como estrategia general.

\subsection{Experimento B: ensamble SARIMAX + XGBoost independiente}\label{sec:por_que_ensemble}

El segundo experimento no está sujeto a un criterio de éxito binario, sino que constituye
un punto de comparación que cualquier evaluación rigurosa de este proyecto debe incluir:
¿existe una combinación más sencilla que los propios componentes de una sola etapa? La
respuesta se construye sin ningún reentrenamiento adicional, combinando de forma puramente
aritmética las predicciones de SARIMAX y de \texttt{xgboost\_alone} ya calculadas en las
secciones anteriores, bajo dos esquemas de ponderación (tabla \ref{tab:ensemble_pesos}):
un promedio simple 50/50 (\texttt{ensemble\_equal}) y un promedio ponderado por el inverso
del MAE de validación de cada componente (\texttt{ensemble\_inverse\_mae}).

\begin{table}[htbp]
  \centering
  \begin{tabular}{lrr}
\toprule
\rowcolor{naranja}
\textbf{Variante} & \textbf{Peso SARIMAX} & \textbf{Peso XGBoost-alone} \\
\midrule
\rowcolor{filaclara} ensemble\_\allowbreak{}equal & 0,5000 & 0,5000 \\
ensemble\_\allowbreak{}inverse\_\allowbreak{}mae & 0,4724 & 0,5276 \\
\bottomrule
\end{tabular}

  \caption{\headlinecolor{Pesos de los dos componentes en cada variante del ensamble.}}
  \label{tab:ensemble_pesos}
\end{table}

Como ya adelantó la tabla \ref{tab:comparacion_test}, ambas variantes no solo son
competitivas, sino que \textbf{superan a los dos componentes individuales}; el mecanismo
de esa ganancia se desarrolla en la sección \ref{sec:mecanismo_ensemble}. Los dos esquemas de ponderación producen resultados casi
idénticos entre sí (una diferencia de MAE de test de apenas el 0,03\,\%), lo que no es
sorprendente dado que los pesos óptimos por inverso del MAE, 0,4724 y 0,5276, se sitúan
muy próximos a un reparto igualitario. Por ello, y pese a que \texttt{ensemble\_inverse\_mae}
obtiene el MAE de test nominalmente más bajo de todo el proyecto, se recomienda
\texttt{ensemble\_equal} como modelo de referencia para la memoria: ofrece un
rendimiento prácticamente indistinguible sin necesitar ningún paso de ajuste que
justificar y, a diferencia de la variante ponderada por inverso del MAE —cuyos pesos se
derivan del propio MAE de validación y hacen, por tanto, que su cifra de validación sea
levemente optimista—, no depende en ningún grado de la partición de validación para
fijar sus coeficientes.

\section{Por qué gana el ensamble: descorrelación de errores}\label{sec:mecanismo_ensemble}

La razón mecánica por la que el ensamble supera a ambos componentes individuales, y no
solo el hecho empírico de que lo hace, se puede hacer explícita observando directamente
la relación entre los errores de SARIMAX y de \texttt{xgboost\_alone} en cada día de la
partición test (figura \ref{fig:dispersion_errores}).

\begin{figure}[htbp]
  \centering
  \includegraphics[width=0.78\textwidth]{fig_dispersion_errores_sarimax_xgb.png}
  \caption{\headlinecolor{Dispersión de residuos SARIMAX frente a XGBoost independiente en la partición test -- evidencia mecánica de la reducción de varianza del ensamble (correlación débil entre ambos errores).}}
  \label{fig:dispersion_errores}
\end{figure}

La correlación entre ambos residuos es de 0,576: positiva, como cabría esperar de dos
modelos que atacan el mismo fenómeno subyacente y comparten, por tanto, una parte de su
señal de error, pero claramente inferior a la unidad. Esa distancia respecto a la
correlación perfecta es, precisamente, la que hace mecánicamente ventajoso promediar los
dos modelos. Para dos pronósticos con varianza de error igual $\sigma^2$ y correlación
$\rho$ entre sus errores, la varianza del error de su media aritmética es
\[
  \operatorname{Var}\!\left(\frac{\hat e_1 + \hat e_2}{2}\right)
    = \frac{\sigma^2 (1 + \rho)}{2},
\]
la expresión clásica de la combinación de pronósticos de
\textcite{bates1969combination}: con $\rho = 1$, errores perfectamente correlacionados,
promediar no reduce la varianza en absoluto, mientras que con $\rho = 0{,}576$ la
varianza del error combinado se reduce a aproximadamente el \textbf{79\,\%} de la varianza
individual, incluso bajo el supuesto simplificador de que ambos componentes tuvieran
exactamente el mismo error cuadrático medio. Esta reducción de varianza por promediado de
predictores parcialmente descorrelacionados es, además, el mismo principio que sustenta el
uso de comités o \emph{ensembles} de modelos en el aprendizaje automático en general
\parencite{perrone1993networks}, aplicado aquí no a una colección de redes con la misma
arquitectura sino a dos familias de modelos deliberadamente distintas.

La tabla \ref{tab:desglose_dia_test} explica además por qué esa correlación es
precisamente 0,576 y no un valor más cercano a la unidad: SARIMAX y
\texttt{xgboost\_alone} no cometen simplemente errores de magnitud distinta sobre los
mismos días, sino errores grandes en \emph{días distintos} (el primero en laborables
ordinarios comparativamente, el segundo en fin de semana y festivos), de modo que su
correlación de error refleja una complementariedad estructural de calendario y no solo
ruido estadístico independiente. La ganancia observada, un MAE de test de 139.681 para
\texttt{ensemble\_inverse\_mae}, un 10,5\,\% por debajo del mejor componente individual
(\texttt{xgboost\_alone}, 156.121), es coherente con el mecanismo de reducción de varianza
descrito por Bates y Granger y con la correlación imperfecta observada entre los errores de
SARIMAX y XGBoost; la identidad, formulada en varianza y bajo el supuesto de varianzas
iguales, no la demuestra, pero la ganancia tiene el signo y el orden de magnitud que ese
mecanismo anticipa.

\section{Anatomía del residuo de la etapa 1}\label{sec:anatomia_residuo_etapa1}

La sección \ref{sec:resultado_negativo_teoria} sostuvo, sobre una base teórica, que el
híbrido pierde porque su etapa 1 es un cuello de botella: una LSTM univariante sobre unas
860 observaciones diarias no puede ver el calendario, y la etapa 2 arranca entonces desde
un residuo cuya desviación típica, 555.385 viajes diarios sobre la serie contigua de
validación y test, es del orden de la propia señal. Esta sección deja de afirmarlo y lo
mide. Es un análisis \textbf{confirmatorio, no exploratorio}: que una red que nunca observa
el calendario deje estructura de calendario en su residuo es casi seguro de antemano; la
aportación es cuantificar \emph{cuánta}, con criterios de aceptación fijados en el código
(constantes de módulo de \texttt{stage1\_residual\_anatomy.py}) antes de observar ningún
resultado. Esta sección y las dos siguientes (\ref{sec:paridad_informativa},
\ref{sec:desglose_subgrupos}) son \textbf{controles diagnósticos}: ninguna de sus cifras
entra en la tabla maestra de comparación, en \texttt{dashboard\_data.MODEL\_NAMES} ni en
\texttt{models/}, y ninguna compite por ser el mejor modelo; una prueba por fase
(\texttt{test\_phase7/8/9\_adds\_no\_model\_to\_the\_master\_comparison}) lo fija
mecánicamente. En esta, además, no se entrena, reajusta ni re-predice ningún modelo, con
la única excepción declarada de la curva de aprendizaje de la sección
\ref{sec:tamano_muestral}.

\subsection{Estructura que retiene el residuo}\label{sec:residuo_estructura}

La tabla \ref{tab:autocorrelacion_residuo} y la figura \ref{fig:acf_pacf_residuo} recogen
la autocorrelación del residuo de la etapa 1 hasta cinco ciclos semanales. La memoria
temporal es \textbf{débil y ambigua}: el mayor valor absoluto de la ACF en los retardos 1
a 35 es de solo 0,166 (retardo 28), y aunque cinco de los treinta y cinco retardos quedan
fuera de la banda de Bartlett, frente a los dos que cabría esperar por azar al 5\,\%,
ninguna autocorrelación alcanza el umbral de 0,20 fijado como evidencia de estructura
persistente. El estadístico de Ljung-Box rechaza la hipótesis de ruido blanco en los tres
retardos examinados ($p < 0{,}01$ en 7, 14 y 28), pero esa prueba es
\textbf{anticonservadora} en este contexto —su corrección de grados de libertad no está
definida para una LSTM sin recuento de parámetros significativo, y sobre un residuo
visiblemente autocorrelacionado rechaza con facilidad—, de modo que ninguna afirmación de
esta sección descansa en su valor $p$: el criterio primario es siempre el tamaño del
efecto.

\begin{figure}[htbp]
  \centering
  \includegraphics[width=0.78\textwidth]{fig_acf_pacf_residuo_lstm.png}
  \caption{\headlinecolor{ACF y PACF del residuo de la etapa 1 hasta el retardo 35, con la banda de Bartlett sombreada y los retardos semanales (7, 14, 21, 28) marcados. El mayor valor absoluto de la ACF es 0,166 en el retardo 28.}}
  \label{fig:acf_pacf_residuo}
\end{figure}

El periodograma (figura \ref{fig:periodograma_residuo}) matiza esta lectura para la
estacionalidad semanal. La cuota de potencia espectral concentrada exactamente en el intervalo espectral (\emph{bin})
de periodo 7 días es solo 1,4 veces la cuota media, por debajo del factor 3 fijado como
umbral, pero esto se debe en buena medida a la fuga espectral: con 393 observaciones no
existe un \emph{bin} de Fourier exactamente en 7,0 días, y la potencia semanal se reparte entre
los \emph{bins} adyacentes (6,4 y 6,7 días) y, sobre todo, hacia el armónico de 7/3 días, que es
el pico individual más alto de todo el periodograma. La ACF en el retardo 7 (0,162) y, muy
especialmente, el perfil por día de la semana apuntan, por tanto, a una estructura semanal
real que la prueba de \emph{bin} único infravalora.

\begin{figure}[htbp]
  \centering
  \includegraphics[width=0.78\textwidth]{fig_periodograma_residuo_lstm.png}
  \caption{\headlinecolor{Periodograma del residuo de la etapa 1: cuota de potencia espectral frente al periodo en días (escala logarítmica). El periodo de 7 días está anotado; la potencia semanal se dispersa entre \emph{bins} adyacentes y el armónico de 7/3 días.}}
  \label{fig:periodograma_residuo}
\end{figure}

La tabla \ref{tab:residuo_dia_semana} y la figura \ref{fig:residuo_dia_semana} muestran ese
perfil por día de la semana. La dispersión entre la media de residuo más alta y la más baja
equivale a 0,60 veces la desviación típica del residuo, muy por encima del umbral de 0,30
$\sigma$ para considerar que hay estructura: la etapa 1 sobrepredice sistemáticamente los
sábados (residuo medio $-190.065$, es decir $-0{,}34\,\sigma$) y los jueves ($-0{,}31\,\sigma$),
e infrapredice los lunes y los viernes ($+0{,}26$ y $+0{,}23\,\sigma$). Como la prueba de
significación entre grupos de día viola la independencia (los residuos están
autocorrelacionados), es este tamaño de efecto, y no un valor $p$, el que sostiene la
conclusión de que persiste un sesgo semanal sistemático que la etapa 1 no elimina.

\begin{figure}[htbp]
  \centering
  \includegraphics[width=0.78\textwidth]{fig_residuo_por_dia_semana.png}
  \caption{\headlinecolor{Residuo medio de la etapa 1 por día de la semana, con barras de $\pm 1$ error típico y el tamaño muestral bajo cada etiqueta. Una barra positiva indica infrapredicción media de ese día.}}
  \label{fig:residuo_dia_semana}
\end{figure}

Donde el residuo sí retiene estructura con claridad es en el calendario irregular
(tabla \ref{tab:residuo_calendario}). El error absoluto medio de la etapa 1 en los días
festivos, 1.818.580 viajes, es \textbf{6,52 veces} el de los días laborables ordinarios
(279.107), y en los días laborables con puente es 2,96 veces, muy por encima del factor 2
fijado como umbral de concentración. Esta es exactamente la señal que la sección
\ref{sec:importancia_variables} identificó del lado de la etapa 2, donde
\texttt{feat\_day\_type\_festivo} es la variable de mayor ganancia del modelo residual, y
coincide con el diagnóstico exploratorio de la fase 4 sobre el periodo completo. Es también
el único de los cuatro criterios de aceptación prerregistrados que cae con claridad en el
lado de ``la estructura persiste'': la etapa 1 univariante deja sin explicar,
precisamente, la irregularidad de calendario que por construcción no puede ver.

\subsection{Señal meteorológica retardada en el residuo}\label{sec:residuo_meteo_anatomia}

Cabría preguntarse si, además del calendario, el residuo de la etapa 1 retiene señal
meteorológica que la etapa 2 no haya sabido extraer. La tabla \ref{tab:residuo_meteo} y la
figura \ref{fig:residuo_meteo} responden que \textbf{no de forma detectable}. De las 105
variables meteorológicas retardadas (lags de 1, 2, 3 y 7 días y media móvil de 7 días; la
meteorología del mismo día está excluida por diseño y la prueba lo reverifica), ninguna
mantiene una correlación de Spearman \emph{parcial} estadísticamente significativa con el
residuo tras controlar por los cuatro términos anuales de Fourier y aplicar la corrección
de Benjamini-Hochberg al 10\,\%. La brecha entre la correlación bruta y la parcial es,
en sí misma, el hallazgo: la humedad relativa máxima retardada un día y la temperatura
máxima retardada dos días presentan correlaciones brutas de $0{,}28$ en valor absoluto,
pero al partializar por el ciclo anual, que ambas comparten con la demanda, colapsan a
$0{,}11$ y $0{,}09$, y no sobreviven a la corrección por comparaciones múltiples.

La lectura correcta de este resultado nulo, a la luz de la objeción metodológica registrada
para esta fase, no es que la arquitectura híbrida debiera haber funcionado: la etapa 2
\emph{recibió} esas 105 variables —de hecho, la sección \ref{sec:importancia_variables}
mostró que la meteorología retardada concentra el 51,6\,\% de su ganancia— y aun así el
híbrido pierde. Lo que este análisis indica es que no se observa señal meteorológica
retardada sin explotar en el residuo; lo que el residuo retiene es estructura de calendario,
no de clima. Que un bloque de variables acapare la mitad de la ganancia por impureza y a
la vez no mantenga ninguna correlación parcial detectable con el residuo es solo
aparentemente contradictorio; la sección \ref{sec:lectura_conjunta} reconcilia ambos
hechos con el tercero que aporta la fase siguiente, la ablación del bloque completo.

\begin{figure}[htbp]
  \centering
  \includegraphics[width=0.78\textwidth]{fig_residuo_vs_meteo.png}
  \caption{\headlinecolor{Correlación de Spearman del residuo de la etapa 1 con las 15 variables meteorológicas retardadas de mayor $|\rho$ parcial$|$. El marcador hueco muestra la $\rho$ bruta: la distancia entre ambos es el efecto del ciclo anual compartido.}}
  \label{fig:residuo_meteo}
\end{figure}

\subsection{La cadena de error, cuantificada}\label{sec:cadena_error}

La tabla \ref{tab:cadena_error} y la figura \ref{fig:cadena_error} cuantifican cuánto de la
distancia entre una LSTM univariante y un XGBoost independiente cierra realmente la etapa
residual. No es una descomposición (\texttt{lstm\_alone}, \texttt{hybrid} y
\texttt{xgboost\_alone} son tres modelos distintos, no sumandos de un mismo error), sino
una cadena de tres puntos evaluados con el mismo \texttt{metrics.all\_metrics} que la tabla
maestra, de modo que no puede contradecirla. Sobre la partición test, la distancia
completa entre \texttt{lstm\_alone} (MAE 280.952) y \texttt{xgboost\_alone} (156.121) es de
124.831 viajes diarios. La etapa 2 recorta 46.729 de esa distancia, el \textbf{37,4\,\%}
del hueco, dejando todavía 78.103 (el 62,6\,\%) de separación entre el híbrido corregido y
un XGBoost independiente sobre las mismas 161 variables. Sobre validación la corrección es
proporcionalmente mucho mayor (130.405, el 78,4\,\% de un hueco más ancho porque
\texttt{lstm\_alone} rinde bastante peor en esa partición), pero los 35.895 viajes
restantes siguen favoreciendo a \texttt{xgboost\_alone}. La conclusión es la misma en ambas
particiones y coincide con lo anticipado en la sección \ref{sec:contraste_central}: la
etapa residual es una corrección real y sustancial, pero parte de tan atrás que la
corrección, por grande que sea, no alcanza a un modelo de una sola etapa.

\begin{figure}[htbp]
  \centering
  \includegraphics[width=0.78\textwidth]{fig_cadena_error.png}
  \caption{\headlinecolor{Barras escalonadas del MAE de test de los tres modelos de la cadena, con el paso entre barras consecutivas anotado en MAE absoluto y como cuota del hueco total. Deliberadamente no es una cascada con conectores.}}
  \label{fig:cadena_error}
\end{figure}

\subsection{Contexto de tamaño muestral}\label{sec:tamano_muestral}

Queda por sustentar con evidencia la afirmación de que la debilidad de la etapa 1 es, en
lo esencial, un problema de tamaño de muestra. Un primer indicio, solo un indicio, procede
de los cinco folds de validación \emph{out-of-fold}. El
MAE por fold no decrece de forma monótona al crecer el tamaño de entrenamiento: cae de
1.461.905 (fold 1, 200 filas) a 282.122 (fold 3, 474 filas) y vuelve a subir a unos
500.000 en los folds 4 y 5, porque cada fold predice además un periodo distinto y de
dificultad distinta. Expresar el error de cada fold como \emph{ratio de habilidad} frente a
\texttt{seasonal\_naive} evaluado sobre exactamente el mismo bloque de fechas atenúa en
parte ese factor de confusión: los folds 3 y 5 alcanzan la paridad con el baseline
estacional (ratios de 0,95 y 0,98), el fold 1 permanece degenerado (2,67) y los folds 2 y
4 quedan por encima de 1. Con cinco puntos, no monótonos y confundidos con la dificultad
del periodo evaluado, esto es contexto indicativo y no una curva de aprendizaje.

La evidencia principal sobre el papel del tamaño de muestra procede de una curva de
aprendizaje propiamente dicha (tabla \ref{tab:curva_aprendizaje} y figura
\ref{fig:curva_aprendizaje}), el único experimento de esta fase que reentrena. Se ajusta
una LSTM de etapa 1, con la arquitectura de la fase 4, sobre ventanas de entrenamiento de
tamaño creciente, en seis fracciones del bloque de entrenamiento post-calentamiento (de
222 a 889 filas), con tres semillas por fracción para poder leer la mediana y la banda
mín-máx en lugar de una única ejecución ruidosa. El conjunto de evaluación se mantiene
fijo (siempre la partición de validación completa) y la partición de test no se consulta
en ningún momento. La ventana de entrenamiento es \textbf{descendente}, anclada en
\texttt{train\_end} y creciendo hacia atrás: esta elección es deliberadamente
conservadora para un experimento de tamaño de muestra, porque las fracciones más pequeñas
entrenan sobre las filas temporalmente más próximas a validación y disfrutan, por tanto, de
una ventaja de recencia; el diseño está sesgado \emph{en contra} de encontrar un efecto de
tamaño de muestra, de modo que un efecto que sobreviva a ese sesgo es más creíble y un
resultado plano es correspondientemente más débil de lo que parecería a primera vista.

El resultado es una historia de tamaño de muestra \textbf{parcial}, y se reporta como tal.
Por un lado, el tamaño de
muestra importa, y mucho, en el extremo bajo: con 222 filas de entrenamiento el modelo
\textbf{colapsa} (las tres semillas producen ajustes degenerados, con un cociente de
desviaciones típicas entre 0,30 y 0,43 y una correlación con la serie real en torno a
0,45, es decir, predicciones cuasi-constantes), exactamente el mismo modo de fallo
documentado para el fold 1 en la sección \ref{sec:esquema_oof}. Pasar de 222 a 356 filas
reduce el MAE de validación mediano de 1.167.187 a 515.603, y a partir de ahí el descenso
continúa, pero se aplana: 502.449 con 489 filas, 490.773 con 622, y ya prácticamente estable
en torno a 411.000-416.000 con 756 y 889 filas, dentro del ruido entre semillas.

Por otro lado, y esto es lo que la curva sí muestra, esa meseta se sitúa muy por encima de ambas
referencias de una sola etapa: el mejor punto LSTM de toda la curva, unos 411.000 de MAE de
validación, sigue estando aproximadamente 166.000 viajes por encima de
\texttt{xgboost\_alone} (244.843) y 138.000 por encima de SARIMAX (273.467), y las tres
fracciones superiores no muestran ninguna tendencia a acercarse a esas líneas. La lectura
correcta, por tanto, no es que ``con más datos la vía LSTM ganaría'': el tamaño de muestra
explica por qué la etapa 1 es \emph{muy} débil en lugar de solo débil —el colapso por
debajo de unas 350 filas y la fuerte ganancia inicial son reales—, pero el techo de una
arquitectura univariante sobre esta serie se encuentra, incluso con todos los datos
disponibles y bajo un diseño de ventana que la favorece, claramente por encima de los
modelos de una sola etapa a los que el híbrido aspira a superar.

\begin{figure}[htbp]
  \centering
  \includegraphics[width=0.78\textwidth]{fig_curva_aprendizaje_lstm.png}
  \caption{\headlinecolor{MAE de validación de la LSTM de etapa 1 frente a las filas de entrenamiento (ventana descendente anclada en \texttt{train\_end}): mediana y banda mín-máx de 3 semillas (5 en la fracción 1,00), con las líneas de referencia de SARIMAX y \texttt{xgboost\_alone} en validación. El marcador hueco indica la fracción cuyas tres semillas degeneraron. La curva se aplana muy por encima de ambas referencias.}}
  \label{fig:curva_aprendizaje}
\end{figure}

\section{Paridad informativa entre etapas: ablación y control LSTM}\label{sec:paridad_informativa}

La sección \ref{sec:anatomia_residuo_etapa1} midió \emph{qué} retiene el residuo de la
etapa 1; esta sección se pregunta algo distinto: si la partición del problema en dos etapas
está \emph{informativamente} justificada. El análisis anterior deja dos cuestiones
abiertas. La primera (Q1): la etapa 1 es univariante mientras que
XGBoost ve 161 variables desde el principio, así que ambas etapas no son equivalentes en
información; ¿pierde la etapa 1 por ser una \emph{LSTM} o por ser \emph{univariante}? La
segunda (Q2): ¿integra ya el XGBoost independiente los efectos temporales y exógenos de forma
simultánea a través de las variables construidas, de modo que separar el problema en dos
etapas añade complejidad sin añadir información predictiva? Esta sección deja de conjeturar
y lo mide con dos controles. Como la sección \ref{sec:anatomia_residuo_etapa1}, es un
análisis \textbf{confirmatorio} y un control diagnóstico bajo el régimen ya declarado allí;
sus dos controles sí reentrenan, pero no persisten ningún artefacto bajo \texttt{models/}.

\subsection{Ablación de bloques de variables sobre \texttt{xgboost\_alone}}\label{sec:ablacion_bloques}

Se reentrena \texttt{xgboost\_alone} sobre subconjuntos restringidos de variables,
manteniendo idénticos la partición de entrenamiento (889 filas tras el descarte de
calentamiento), la rejilla de 54 puntos, el \emph{holdout} cronológico interno y el
protocolo de evaluación; lo único que cambia es qué bloques son visibles. Los subconjuntos
—\texttt{completo} (161 variables), \texttt{solo\_temporal} (39: retardos de \texttt{total}
y de operadores más medias móviles), \texttt{solo\_exogeno} (122: calendario, Fourier y
meteorología), \texttt{sin\_meteo} (56) y \texttt{solo\_calendario} (17)— se fijaron
\emph{antes} de ver ningún resultado, y el criterio de materialidad (un bloque aporta si
retirarlo eleva el MAE de validación en al menos un 5\,\% respecto a \texttt{completo}) es
una constante de módulo evaluada en código, no mediante inspección manual. El umbral del 5\,\% absorbe el ruido
de rejilla y semilla, ya que cada subconjunto rehace su propia búsqueda de
hiperparámetros.

Esto es un \textbf{estudio de retirada de bloques, no una descomposición}. Los bloques
\textbf{no son disjuntos en información}: \texttt{fourier\_annual} y la meteorología
retardada codifican ambos el ciclo anual —la sección \ref{sec:residuo_meteo_anatomia} lo
midió, viendo colapsar las correlaciones al parcializar por los términos anuales—, y los
retardos 7/14/21/28 de \texttt{total} codifican implícitamente el calendario semanal. Por
tanto, \texttt{solo\_temporal} sigue ``sabiendo'' el día de la semana y \texttt{solo\_exogeno}
sigue ``sabiendo'' la estación. Las columnas de la tabla \ref{tab:ablacion_bloques} no
suman a nada y no deben leerse como contribuciones aditivas.

El resultado responde a Q2 de la forma prevista: \textbf{\texttt{solo\_temporal} y
\texttt{solo\_exogeno} empeoran ambos por encima del umbral}, un 88,6\,\% y un 30,9\,\% de
MAE de validación respectivamente, de modo que el modelo de una sola etapa
\textbf{integra internamente los dos tipos de información}. Esto es la corroboración
mecánica del resultado ya establecido en la sección \ref{sec:contraste_central} (no una
prueba nueva; el capítulo ya mostró la mitad predictiva, que el híbrido pierde): la
división en dos etapas hace, a través de etapas, lo que un único modelo tabular ya hace en
una. Dos lecturas secundarias, reportadas sin veredicto de aprobado/suspenso:
\texttt{solo\_calendario} (17 variables) queda un 22,8\,\% por encima de \texttt{completo}
, sustancial pero lejos de la paridad, así que la serie no está determinada solo por el
calendario determinista; y \texttt{sin\_meteo} rinde un 3,2\,\% \emph{mejor} que
\texttt{completo}, es decir, el bloque de 105 columnas meteorológicas no se gana su sitio y
actúa como ruido neto leve, en línea con el hallazgo de la sección
\ref{sec:residuo_meteo_anatomia} de que la meteorología retardada es cuasi-ruido-blanco en
el residuo. La figura \ref{fig:ablacion_bloques} ordena los cinco subconjuntos por MAE de
validación y superpone el MAE de test de cada uno, de modo que el orden de la ablación
pueda leerse directamente junto con su estabilidad entre particiones.

\begin{figure}[htbp]
  \centering
  \includegraphics[width=0.78\textwidth]{fig_ablacion_bloques.png}
  \caption{\headlinecolor{MAE de \texttt{xgboost\_alone} reentrenado sobre subconjuntos restringidos de variables, mejor subconjunto arriba. La línea discontinua marca el MAE de validación de \texttt{completo} (244.843) y los rombos el MAE de test como comprobación de estabilidad. Retirar el bloque temporal o el exógeno empeora el modelo por encima del umbral del 5\,\%; retirar la meteorología no.}}
  \label{fig:ablacion_bloques}
\end{figure}

\subsection{LSTM en paridad informativa}\label{sec:paridad_lstm}

Para responder a Q1 se dota a la rama recurrente de \textbf{paridad informativa}: la ventana
de 28 pasos del objetivo alimenta la LSTM y el vector de variables predictoras \texttt{feat\_}
del \emph{propio día $t$} se concatena después
(\texttt{Input}$(W,1)\to$\texttt{LSTM}$\to$\texttt{concat}(\texttt{Input}$(k)$)$\to$
\texttt{Dropout}$\to$\texttt{Dense}$(1)$). Es el cambio mínimo que logra paridad real: una
ventana multivariante $(W,k)$ daría el calendario de los 28 días previos, pero no el del día
que se predice, porque la ventana abarca $[t-W, t)$ y excluye la fila $t$ por construcción.
Cada columna es la misma fila, sujeta a las mismas guardias anti-fuga, que ve XGBoost
(guardas de las fases 2 y 3: sin operadores del mismo día, sin meteorología del mismo
día), y una aserción
reverifica el prefijo \texttt{feat\_} en el punto de entrada.

La interpretación es asimétrica y se declara antes de los resultados (objeción O4):
la LSTM en paridad ve una ventana cruda de 28 pasos \emph{más} el vector del día $t$, es
decir, estrictamente \emph{más} información que el vector plano de XGBoost (que solo dispone
de 7 retardos seleccionados y 4 estadísticos móviles). En consecuencia, \textbf{una derrota
clara sería una conclusión fuerte} (la asimetría informativa quedaría descartada como causa
dominante) y \textbf{una victoria es débil} (procede en parte de información extra, no solo
de la arquitectura). Y un buen resultado aquí sería un argumento \emph{contra} el híbrido,
no a favor de uno mejor: una etapa 1 que ya ve la matriz de la etapa 2 no deja a esta nada
que corregir que la etapa 1 no pudiera ver.

\textbf{Base de comparación y número de semillas.} El punto de referencia es la LSTM
univariante de la fase 4. Se compara contra su mediana de cinco semillas
(42, 1, 2, 3, 4) en la fracción 1,00 de la curva de aprendizaje de la sección
\ref{sec:tamano_muestral}, \textbf{416.324} de MAE de validación, y no contra su mejor
semilla individual (411.142; cuadro \ref{tab:paridad_lstm}). El veredicto de paridad se lee sobre una mediana de
cinco semillas, así que la base tiene que ser también una mediana de cinco semillas (mismo
arnés, mismas semillas, misma partición de validación, genuinamente comparable); enfrentar
una mediana contra un mejor caso sesgaría el cierre del hueco a la baja, hacia la
conclusión que el TFM ya defiende, que es justo donde el análisis debe ser más estricto. La
mediana de cinco semillas resulta coincidir con la de tres (la semilla 1 es el valor
central en ambos casos), de modo que el hueco a cerrar sigue siendo $416.324 - 244.843 =
171.481$ y los umbrales no cambian.

El paso de tres a cinco semillas se decidió por la
regla prerregistrada (secciones de método y objeción O6: subir a cinco solo si la sonda de
tiempos muestra holgura, decidido \emph{antes} de la ejecución); la sonda mostró amplia
holgura (el ajuste más lento fue de 9,7\,s frente a un techo de unos 6\,min). La sonda
también dejó visibles unos MAE de validación provisionales con semilla 42 (395.274 en
\texttt{calendario}, 369.835 en \texttt{completo}), ambos dentro de la banda ambigua del
25--60\,\%; se hace constar explícitamente que esas cifras \emph{eran} visibles antes de
elegir el número de semillas y que la regla, basada en holgura de cómputo, gobierna con
independencia de dónde cayeran.

\textbf{Lecturas prerregistradas.} Con un hueco de 171.481, cerrar el $\geq$ 60\,\%
(MAE de validación mediano $\leq 313.435$) se lee como \emph{asimetría informativa}: la
debilidad de la etapa 1 está dominada por lo que no se le dejó ver. Cerrar el $\leq$ 25\,\%
(mediano $\geq 373.454$) se lee como \emph{limitación arquitectónica}: un modelo
recurrente simplemente no es competitivo en este problema a este tamaño de muestra, ni
siquiera en paridad informativa. Entre ambos, \emph{mixto/ambiguo}. La degeneración se
evalúa primero: si la mayoría de semillas de un ámbito degenera, el veredicto es
\emph{fallo de capacidad/muestra}. Y si la banda mín-máx cruza un umbral, el veredicto se
degrada a \emph{mixto/ambiguo} con independencia de dónde caiga la mediana (objeción O6).

\textbf{Resultado.} Ninguno de los dos ámbitos se acerca a cerrar el hueco (tabla
\ref{tab:paridad_lstm}, figura \ref{fig:paridad_lstm}). Ninguna semilla degenera y ninguna
queda infra-entrenada (\texttt{best\_epoch} mediano 75 y 53, muy por encima del umbral de
20). El ámbito \texttt{calendario} ($k=17$), el más nítido para la afirmación de ceguera al
calendario, cierra solo un \textbf{16,2\,\%} del hueco (MAE de validación mediano 388.625;
banda 382.695--395.877, sin cruce de umbral): veredicto \emph{limitación
arquitectónica}. Es decir, dada exactamente la información de calendario y Fourier donde la
sección \ref{sec:residuo_estructura} localizó la estructura del residuo, un modelo
recurrente en paridad informativa apenas recupera una sexta parte de la distancia hasta
XGBoost. El ámbito \texttt{completo} ($k=161$) lo hace mejor, cierra un \textbf{31,0\,\%}
(mediano 363.133), pero su banda (354.235--391.630) cruza el umbral del 25\,\%, de modo que
el veredicto se degrada a \emph{mixto/ambiguo}; y por la asimetría O4 un cierre parcial es
de todos modos evidencia débil. En conjunto, la respuesta a Q1 es que la etapa 1 pierde
principalmente por ser \emph{recurrente sobre esta serie y a este tamaño de muestra}, no por
ser univariante: la paridad informativa no la rescata, aunque el veredicto
\emph{mixto/ambiguo} del ámbito \texttt{completo} impide cerrar del todo la explicación
informativa.

\begin{figure}[htbp]
  \centering
  \includegraphics[width=0.78\textwidth]{fig_paridad_lstm.png}
  \caption{\headlinecolor{MAE de validación de la LSTM en paridad informativa por ámbito: mediana y banda mín-máx de 5 semillas, con las líneas de referencia de la LSTM univariante (mediana de 5 semillas), SARIMAX y \texttt{xgboost\_alone} en validación. Ninguno de los dos ámbitos se acerca a la línea de \texttt{xgboost\_alone}.}}
  \label{fig:paridad_lstm}
\end{figure}

\subsection{Lectura conjunta}\label{sec:lectura_conjunta}

Los dos controles apuntan en la misma dirección. La ablación muestra que
\texttt{xgboost\_alone} ya usa de forma medible tanto la información temporal como la
exógena: no hay un componente redundante que la etapa recurrente del híbrido pudiera
aportar por separado. El control de paridad muestra que, aun entregándole a una LSTM
exactamente las variables predictoras que ve XGBoost, e incluso más información en forma de
ventana cruda, la vía recurrente no alcanza a los modelos de una sola etapa. La debilidad
de la etapa 1, por tanto, no parece explicarse en lo esencial por el acceso a la
información: la lectura más plausible es la combinación de una arquitectura recurrente y
unas 860 observaciones diarias, con la salvedad, que el ámbito \texttt{completo} deja
abierta, de que una fracción del hueco pueda deberse también al acceso o a la
representación de la información (sección \ref{sec:limitaciones}).
Esto refuerza, por una vía independiente, la conclusión del contraste central
(\ref{sec:contraste_central}): para la demanda agregada diaria a este tamaño de muestra, un
único modelo de \emph{gradient boosting} sobre una matriz de variables predictoras bien
construida supera a
un híbrido residual secuencial, y el beneficio del híbrido se limita a compensar una etapa
temporal débil en lugar de a batir a modelos de una sola etapa fuertes.

La ablación permite además cerrar una tensión que arrastran las secciones anteriores.
Tres hechos sobre la meteorología retardada parecen
incompatibles: la sección \ref{sec:importancia_variables} reporta que ese bloque concentra
el 51,6\,\% de la ganancia del modelo residual; la sección
\ref{sec:residuo_meteo_anatomia} no encuentra ninguna de sus 105 variables con correlación
parcial significativa con el residuo; y la ablación de la sección
\ref{sec:ablacion_bloques} muestra que retirarlo por completo \emph{mejora} el MAE de
validación en un 3,2\,\% (237.040 frente a 244.843).

La reconciliación está en qué mide
realmente la ganancia por impureza. Esa métrica acumula la reducción de impureza sobre
\emph{cada corte} en que interviene una variable y, por tanto, premia sistemáticamente a
las variables continuas de alta cardinalidad: una columna meteorológica continua ofrece al
algoritmo centenares de puntos de corte candidatos, mientras que una indicadora binaria de
calendario ofrece exactamente uno. Con 105 columnas continuas frente a 17 indicadoras y
solo 552 filas de entrenamiento OOF, el bloque meteorológico dispone de una enorme
superficie de corte sobre la que ajustar ruido, y acumula cuota de ganancia sin aportar
señal explotable. Leídos así, los tres hechos dejan de estar en conflicto y el 51,6\,\% no
es una anomalía, sino \textbf{evidencia adicional a favor de la tesis de este capítulo}: la
etapa 2 dedicó la mayor parte de su capacidad de modelado a un bloque de variables que dos
pruebas independientes (una correlacional, otra de retirada) muestran que no aporta nada.
Es coherente con ello que la variable individual de mayor ganancia del modelo residual siga
siendo \texttt{feat\_day\_type\_festivo} (0,125, sección \ref{sec:importancia_variables}):
el mismo fenómeno visto por el otro lado, con la única variable genuinamente informativa
siendo una indicadora de baja cardinalidad cuya cuota \emph{de grupo} queda diluida. La
conclusión metodológica, aplicable más allá de este TFM, es que una cuota de ganancia por
impureza no es evidencia de utilidad predictiva y no debe leerse como tal sin un
contraste de retirada que la respalde.

\section{Desglose por subgrupos y periodos}\label{sec:desglose_subgrupos}

El contraste central deja abierta una pregunta concreta: aunque el híbrido no gane globalmente, podría
comportarse de forma distinta en festivos, puentes, fines de semana o periodos de demanda
anómala. Esta sección la responde de forma sistemática. Bajo el mismo régimen de
control diagnóstico declarado en la sección \ref{sec:anatomia_residuo_etapa1}, es un
reanálisis puro de \texttt{data/\allowbreak{}processed/\allowbreak{}*\_predictions.parquet}:
no se entrena, reajusta ni vuelve a predecir nada, y todas las cifras proceden de
\texttt{metrics.all\_metrics}, la misma ruta que el resto del capítulo.

El riesgo dominante aquí no es la incompletitud, es la búsqueda selectiva de significación (\emph{p-hacking}): buscar subgrupos
hasta encontrar uno que favorezca al híbrido es \emph{p-hacking} con independencia de por
qué se emprenda la búsqueda. Por eso el diseño se fijó \textbf{antes} de calcular ningún MAE, clasificación
o diferencia —solo se contaron los tamaños muestrales—, de modo que ``el híbrido no gana en
ningún subgrupo'' tenga el mismo estatuto que cualquier hallazgo positivo.

\subsection{Diseño prerregistrado}\label{sec:desglose_diseno}

Se fijaron \textbf{quince subgrupos} como constantes de módulo en
\texttt{src/\allowbreak{}evaluation/\allowbreak{}subgroup\_breakdown.py}, construidos a partir de columnas ya
presentes en \texttt{features\_daily.parquet} (\texttt{day\_type}, \texttt{feat\_is\_bridge\_day},
\texttt{feat\_is\_weekend}, \texttt{holiday\_name}): regímenes de calendario (laborable
ordinario, laborable en puente, sábado, domingo, festivo y sus uniones), estación (julio y
agosto frente al resto del año), dos ventanas ancladas en festivo (22 dic--6 ene; Jueves
Santo $\pm$4 días) y demanda anómala. Un día es \emph{anómalo} si
$\lvert \text{total}_t - \text{total}_{t-7}\rvert$ supera $3{,}0$ desviaciones robustas,
con la escala (MAD) estimada \textbf{solo sobre train} para no ajustar la definición a los
datos de evaluación. La regla alternativa de mediana móvil a 28 días se prototipó sobre
train y se descartó: en ella domina la estacionalidad semanal, con lo que sería un
detector de fines de semana que duplicaría en silencio el subgrupo \texttt{domingo}. Esa
variante no se reporta.

El umbral de inferencia es $n \geq 20$ (\texttt{MIN\_N\_INFERENTIAL}), justificado por el
apalancamiento de un solo día: a $n=20$ un día pesa $\leq$5\,\% de la media del grupo; a
$n=5$, un 20\,\%; a $n=3$, un 33\,\%. Cualquier subgrupo por debajo de 20 se calcula y se
tabula con su $n$, se marca ``no inferencial'' y queda excluido de toda afirmación de
victoria o derrota.

\textbf{Test es la partición de veredicto; validación es el requisito de réplica.} Un
subgrupo cuenta como victoria del híbrido solo si este queda primero en \textbf{ambas}
particiones. Este es el dispositivo anti-\emph{p-hacking} más fuerte del diseño: convierte
una búsqueda de nueve subgrupos en una que exige réplica independiente, y no cuesta nada
porque las dos particiones ya están publicadas. La partición combinada \texttt{val\_test}
(393 días contiguos) es un ámbito \textbf{secundario} prerregistrado, usado solo de forma
descriptiva para los subgrupos que no alcanzan $n=20$ en test; su coste es que el orden
SARIMAX, los hiperparámetros de XGBoost y los pesos del ensamble se eligieron sobre
validación, así que sus resultados se etiquetan como secundarios y nunca encabezan una
conclusión.

\textbf{Resolución del ámbito de veredicto para \texttt{demanda\_anomala}.} Este subgrupo
, el único fijado por nombre en el diseño prerregistrado, tiene $n=16$ en test y $n=34$ en
validación, y su vista combinada \texttt{val\_test} alcanza $n=50$. Recibe el veredicto
\emph{no inferencial} en el ámbito de veredicto (test), exactamente igual que cualquier
otro subgrupo por debajo del umbral, \textbf{sin excepción}. Su resultado combinado se
reporta como ámbito secundario y descriptivo, con la salvedad completa ya indicada, y
nunca porta un veredicto ni un titular. El ámbito combinado estaba disponible y
\textbf{deliberadamente no} se promovió a ámbito de veredicto para este subgrupo: elegir
por subgrupo qué partición decide es un grado de libertad del investigador, y adoptarlo
precisamente para el subgrupo fijado de antemano es donde menos defendible
sería. La función de veredicto lee la constante \texttt{VERDICT\_SCOPE} de forma uniforme
para los quince subgrupos y ningún camino de código asigna un ámbito de veredicto por
subgrupo; una prueba automatizada lo fija.

\textbf{Navidad y Semana Santa no admiten conclusión.} La ventana de evaluación de 393
días contiene exactamente una Navidad (solo en validación, $n=16$; test $n=0$) y una
Semana Santa (solo en test, $n=10$; validación $n=0$). Ninguna definición arregla esto: es
una propiedad de la partición, no del umbral. Ambos periodos se reportan de forma
descriptiva, se marcan \emph{no analizable} y no se emite ninguna afirmación de victoria o
derrota sobre Navidad o Semana Santa en este TFM. La lista de subgrupos se fijó por
adelantado y añadir uno tras ver los resultados invalidaría la corrección por comparaciones
múltiples: es la restricción que prevalece sobre la exhaustividad.

\subsection{Resultados por subgrupo}\label{sec:desglose_resultados}

Las tablas \ref{tab:subgrupos_test} y \ref{tab:subgrupos_val} recogen, por subgrupo y para
un conjunto reducido de modelos, el mejor competidor y su MAE, el mejor híbrido y su MAE,
la diferencia entre ambos y el veredicto prerregistrado. La figura
\ref{fig:subgrupos_mae} lo representa con barras huecas para los subgrupos no
inferenciales.

\begin{figure}[htbp]
  \centering
  \includegraphics[width=0.78\textwidth]{fig_subgrupos_mae.png}
  \caption{\headlinecolor{MAE por subgrupo en la partición test para el conjunto reducido de modelos. Las barras huecas y traslúcidas corresponden a subgrupos no inferenciales ($n < 20$).}}
  \label{fig:subgrupos_mae}
\end{figure}

El resultado es negativo en todos los subgrupos inferenciales. De los quince subgrupos prerregistrados, nueve alcanzan $n \geq
20$ en test; cuatro quedan por debajo (laborable en puente $n=3$, festivo $n=5$, calendario
irregular $n=8$, demanda anómala $n=16$) y dos están excluidos por falta de solape
val/test. \textbf{En 0 de los 9 subgrupos con $n$ suficiente gana el híbrido con
evidencia}, y en 0 queda siquiera primero sin evidencia: en los nueve, el mejor híbrido es
peor que el mejor competidor por entre 53.019 (verano) y 133.578 (domingo) viajes diarios
de MAE. En validación el único subgrupo donde un híbrido queda primero es festivo ($n=9$),
que es no inferencial y además no replica el hecho en test. La búsqueda se ha completado con
resultado negativo.

El único hallazgo positivo ya conocido se mantiene y se reitera: en festivos, la etapa 2
recorta el error de la etapa 1 de 1.585.388 a 666.572 (un $-58$\,\%). El mecanismo de
corrección residual es real; lo que es insuficiente es la etapa 1 de la que parte, y ese
recorte no convierte a festivo en una victoria de subgrupo —sigue siendo peor que
\texttt{xgboost\_alone} y que el ensamble, y su $n=5$ lo hace indicativo, no concluyente—.

La etiqueta \texttt{demanda\_anomala} usa el $\text{total}_t$ observado: es legítima para un
reporte a posteriori, pero \textbf{ilegítima como regla de enrutamiento}. No cabe concluir
de ello que el híbrido deba desplegarse en los días anómalos: la anomalía no se conoce
antes de observar la demanda del propio día. Su vista combinada \texttt{val\_test}
($n=50$, secundaria y descriptiva) tampoco cambia el signo: el mejor híbrido es peor que
\texttt{xgboost\_alone} por unos 123.000 de MAE.

\subsection{Incertidumbre y comparaciones múltiples}\label{sec:desglose_incertidumbre}

La incertidumbre de cada diferencia se obtiene con un remuestreo \emph{bootstrap} pareado de bloques
móviles. Como un subgrupo categórico está disperso por la partición, el remuestreo se hace
sobre la serie de evaluación completa y contigua, en bloques de 7 días (preservan la
dependencia semanal, que es la autocorrelación dominante y la razón por la que un contraste
pareado ingenuo sobreestima la significación), con 2000 réplicas y semilla global. La
pertenencia al subgrupo viaja con el día, y \texttt{best\_model} y \texttt{best\_hybrid} se
\textbf{vuelven a seleccionar dentro de cada réplica} (\emph{argmin} sobre competidores y
sobre \{\texttt{hybrid}, \texttt{hybrid\_weighted}\}), de modo que la doble selección se
paga en lugar de ignorarse. Una réplica que aporta menos de 20 días del subgrupo se
descarta; si la fracción de réplicas utilizables baja de $0{,}90$, el intervalo de
confianza se suprime. El cuadro \ref{tab:subgrupos_inferencia} recoge la diferencia, su
intervalo y las dos p, bruta y ajustada, para cada subgrupo de la familia inferencial, y
la figura \ref{fig:subgrupos_delta_ic} representa esas mismas diferencias como gráfico de
bosque frente a la línea del cero.

\begin{figure}[htbp]
  \centering
  \includegraphics[width=0.78\textwidth]{fig_subgrupos_delta_ic.png}
  \caption{\headlinecolor{Gráfico de bosque de la diferencia de MAE por subgrupo con su IC \emph{bootstrap} del 95\,\%. La línea del cero es la referencia: un IC íntegramente a su derecha es evidencia de que el híbrido pierde en ese subgrupo. Marcadores huecos y sin IC para los subgrupos no inferenciales.}}
  \label{fig:subgrupos_delta_ic}
\end{figure}

La corrección de Benjamini-Hochberg se aplica con la misma llamada y el mismo $q=0{,}10$
que el análisis meteorológico de la sección \ref{sec:anatomia_residuo_etapa1}, sobre la
familia de los nueve subgrupos inferenciales (una hipótesis por subgrupo, no una por
subgrupo$\times$modelo, porque el estadístico focal es una única diferencia). Los subgrupos
se solapan fuertemente (\texttt{entre\_semana} contiene a \texttt{laborable\_ordinario}),
así que los contrastes son positivamente dependientes —el régimen (PRDS) bajo el que
\texttt{fdr\_bh} sigue siendo válido—; las diferencias entre subgrupos \textbf{no
descomponen}: es una familia de vistas solapadas, no una partición.

En los nueve subgrupos inferenciales la diferencia es positiva y significativa tras
Benjamini-Hochberg: la evidencia es que el híbrido es significativamente \textbf{peor}, no
que empate. Este es el desenlace más expuesto a sobrelectura, así que su tratamiento —la
corrección sobre la familia inferencial y el requisito de réplica en validación— está
prerregistrado. Donde la inferencia es
imposible —festivo ($n=5$), laborable en puente ($n=3$)— no se realiza inferencia
estadística debido al reducido tamaño muestral: se reporta únicamente el MAE puntual con
su $n$.

\subsection{Vista temporal}\label{sec:desglose_temporal}

La figura \ref{fig:mae_movil_periodo} traza, sobre los 393 días contiguos de validación y
test, el MAE móvil a 28 días de cada modelo (panel superior) y la diferencia móvil entre el
mejor híbrido y el mejor competidor (panel inferior), con la frontera val/test marcada. La
ventana de 28 días se hereda de $W=28$ de la LSTM y de la ventana de las variables móviles;
no se ha ajustado. La diferencia móvil se mantiene por encima de cero en toda la serie: la
posición relativa del híbrido no se invierte en ningún tramo. Es una figura descriptiva:
ventanas consecutivas se solapan en 27 días, así que no se le aplica ningún contraste y las
ventanas que cruzan la frontera mezclan las dos particiones.

\begin{figure}[htbp]
  \centering
  \includegraphics[width=0.78\textwidth]{fig_mae_movil_periodo.png}
  \caption{\headlinecolor{MAE móvil a 28 días por modelo (panel superior) y diferencia móvil mejor híbrido $-$ mejor competidor (panel inferior) sobre la ventana contigua de validación y test. Regla vertical en la frontera val/test. Figura descriptiva: sin contraste, por el fuerte solapamiento de ventanas consecutivas.}}
  \label{fig:mae_movil_periodo}
\end{figure}

\subsection{Lectura}\label{sec:desglose_lectura}

La búsqueda por subgrupos y periodos se ha completado con un
resultado negativo bien instrumentado. En los nueve subgrupos con tamaño muestral
suficiente el híbrido \textbf{no queda primero en ninguno}, ni con evidencia ni sin ella, y en los
nueve la evidencia \emph{bootstrap}, corregida por comparaciones múltiples, es que es
significativamente peor que el mejor modelo de una sola etapa o el ensamble. En los cuatro
subgrupos no inferenciales, incluido el de demanda anómala, no se emite veredicto. El único comportamiento diferencial real del híbrido, el recorte del
$58$\,\% del error de la etapa 1 en festivos, es un mecanismo de corrección, no una
victoria, y ya estaba documentado. Esta sección no abre preguntas de investigación nuevas:
cierra la que quedaba, confirmando por una vía sistemática la respuesta del contraste
central (\ref{sec:contraste_central}).

\section{Respuesta a las preguntas de investigación}\label{sec:respuesta_preguntas}

Se retoman aquí, en el mismo orden en que se plantearon en la sección
\ref{sec:preguntas_investigacion}, las cinco preguntas de investigación que estructuran
este TFM, respondidas de forma explícita y remitiendo a la evidencia ya presentada en
este capítulo.

\begin{description}[nosep,leftmargin=0pt,labelindent=0pt,style=nextline,font=\bfseries]
  \item[PI1. ¿Bate el híbrido a SARIMAX en validación y en test simultáneamente?]
    No, en ninguna de las dos particiones: el MAE del híbrido es 7.270 viajes diarios peor
    en validación y 70.597 peor en test (tabla \ref{tab:delta_hibrido}, sección
    \ref{sec:contraste_central}).

  \item[PI2. ¿Bate el híbrido a \texttt{xgboost\_alone}, entrenado sobre las mismas
    variables predictoras?]
    No, y por un margen todavía mayor que frente a SARIMAX: 35.895 en validación y 78.103
    en test (tabla \ref{tab:delta_hibrido}).

  \item[PI3. ¿Reduce el híbrido específicamente el error en festivos y puentes frente a
    SARIMAX y \texttt{xgboost\_alone}?]
    No: es peor que ambos en prácticamente todos los grupos de calendario de la tabla
    \ref{tab:desglose_dia_test}, con la salvedad de que la comparación en festivos (n=5) y
    puentes (n=3) en la partición test debe leerse como indicativa, no concluyente, dado su
    reducido tamaño muestral (sección \ref{sec:error_dia}).

  \item[PI4. ¿Funciona la etapa 2 como mecanismo de corrección residual, aunque el
    híbrido completo no supere a los modelos de una sola etapa?]
    Sí: reduce el MAE de la etapa 1 en un 16,6\,\% en test y, de forma más marcada todavía,
    en un 58\,\% en festivos, de 1.585.388 a 666.572 (sección \ref{sec:error_dia}), y la
    variable de mayor ganancia del modelo residual, \texttt{feat\_day\_type\_festivo}
    (0,125), apunta, a nivel de variable individual, a que esa corrección procede
    de la señal de calendario que la etapa 1 no puede ver (sección
    \ref{sec:importancia_variables}). El mecanismo de la arquitectura híbrida es real; lo
    que resulta insuficiente es la etapa 1 de la que parte.

  \item[PI5. ¿Existe una combinación más simple que supere a ambos componentes de una
    sola etapa?]
    Sí: el ensamble equiponderado 50/50 de SARIMAX y \texttt{xgboost\_alone}
    (\texttt{ensemble\_equal}), sin ninguna red LSTM ni reentrenamiento adicional, alcanza
    un MAE de test de 139.725; su variante ponderada por el inverso del MAE de validación
    (\texttt{ensemble\_inverse\_mae}, pesos 0,4724/0,5276) lo reduce ligeramente a 139.681,
    un 10,5\,\% por debajo del mejor de sus dos componentes, en consonancia con la
    descorrelación de errores entre ambos (sección \ref{sec:mecanismo_ensemble}).
\end{description}

Los dos controles diagnósticos de la sección \ref{sec:paridad_informativa} no abren
preguntas de investigación nuevas, pero refuerzan las respuestas a PI1, PI2 y PI4 por una
vía independiente: la ablación de bloques (\ref{sec:ablacion_bloques}) confirma
mecánicamente que \texttt{xgboost\_alone} ya integra la información temporal y la exógena
—lo que explica por qué separar el problema en dos etapas no añade poder predictivo—, y el
control de paridad informativa (\ref{sec:paridad_lstm}) reduce la plausibilidad de que la
desventaja de la etapa 1 se explique únicamente por su carácter univariante: sigue por
detrás de los modelos de una sola etapa incluso viendo las mismas variables predictoras que
XGBoost (veredicto \emph{limitación arquitectónica} en \texttt{calendario}). No obstante,
el ámbito \texttt{completo} presenta un resultado \emph{mixto/ambiguo}, por lo que no puede
descartarse que parte de la diferencia esté relacionada también con el acceso o la
representación de la información (sección \ref{sec:limitaciones}).
